% =============================================================================
% test-abnt-unb-B.tex
% Teste visual — abnt-unb-B.cls (Modelo UnB — clássico)
%
% Compilação (da raiz do repositório):
%   TEXINPUTS=./src:./assets/unb: latexmk -pdf -outdir=tests tests/test-abnt-unb-B.tex
%
% Checklist de verificação visual:
%   1  — Logo UnB pequeno centralizado no topo (~2.5cm)
%   2  — Hierarquia institucional uppercase (universidade, instituto, programa)
%   3  — Autor uppercase centralizado
%   4  — Título bold uppercase centralizado
%   5  — Tipo de trabalho centralizado
%   6  — Local uppercase e ano no rodapé
%   7  — Sem faixa azul (diferente do modelo A)
%   8  — Folha de rosto com cabeçalho discreto + natureza do trabalho
%   9  — Elementos herdados funcionam (aprovação, dedicatória, etc.)
%  10  — Compila com pdfLaTeX
% =============================================================================

\documentclass{abnt-unb-B}
\usepackage{abnt-resumo}

% =============================================================================
% METADADOS UnB
% =============================================================================
\faculdade{Instituto de Ciência Política}
\curso{Ciência Política}
\grau{Mestre}
\tipotrabalho{Dissertação de Mestrado}

% =============================================================================
% METADADOS ABNT (herdados de abnt.cls)
% =============================================================================
\titulo{Participação política e redes sociais}
\subtitulo{um estudo comparado entre Brasil e Portugal}
\autor{Ciclana Menezes dos Santos}
\orientador{Prof.\ Dr.\ Beltrano Vieira}
\ano{2024}

\natureza{Dissertação apresentada como requisito parcial para obtenção
do grau de Mestre em Ciência Política pela Universidade de Brasília.}
\programa{Programa de Pós-Graduação em Ciência Política}
\area{Política Comparada}

% Banca examinadora
\dataaprovacao{15 de agosto de 2024}
\membrodabanca{Prof.\ Dr.\ Beltrano Vieira}
  {Universidade de Brasília (Orientador)}
\membrodabanca{Profa.\ Dra.\ Sicrana Ferreira}
  {Universidade de Brasília}
\membrodabanca{Prof.\ Dr.\ Fulano Almeida}
  {Universidade Federal de Minas Gerais}

% =============================================================================
% DOCUMENTO
% =============================================================================
\begin{document}

\pretextual

\imprimircapa
\imprimirfolhaderosto
\imprimirfolhadeaprovacao

\imprimirdedicatoria{%
À memória da minha avó, que sempre acreditou na força da educação.}

\imprimiragradecimentos{%
Agradeço ao meu orientador, Prof.\ Dr.\ Beltrano Vieira, pela
dedicação e pelos ensinamentos ao longo desta pesquisa.

Agradeço à CAPES pelo financiamento desta pesquisa
(Portaria 206/2018).}

\imprimirepigrafe{%
\textit{``A democracia é o governo do povo, pelo povo, para o povo.''}
--- Abraham Lincoln}

\imprimirresumo{%
Esta dissertação investiga a relação entre participação política e
uso de redes sociais em perspectiva comparada entre Brasil e Portugal.
A partir de dados de survey aplicados em ambos os países no período
de 2020 a 2023, analisa-se como diferentes plataformas digitais
influenciam o engajamento cívico dos cidadãos. Os resultados indicam
que, embora as redes sociais ampliem o acesso à informação política,
seu impacto na participação efetiva varia significativamente conforme
o contexto institucional e a cultura política de cada país.%
}{participação política; redes sociais; política comparada;
democracia digital; Brasil; Portugal}

\imprimirabstract{%
This dissertation investigates the relationship between political
participation and social media use in a comparative perspective
between Brazil and Portugal. Based on survey data collected in both
countries from 2020 to 2023, it analyzes how different digital
platforms influence citizens' civic engagement. Results indicate
that, although social media broaden access to political information,
their impact on effective participation varies significantly
according to each country's institutional context and political
culture.%
}{political participation; social media; comparative politics;
digital democracy; Brazil; Portugal}

\imprimirsumario

% --- Textuais ---
\textual

\chapter{Introdução}

A participação política constitui um dos pilares fundamentais das
democracias contemporâneas. Com a expansão das tecnologias digitais,
novas formas de engajamento cívico emergiram, alterando a dinâmica
entre cidadãos e instituições políticas.

\section{Problema de pesquisa}

De que maneira o uso de redes sociais influencia a participação
política em democracias com diferentes tradições institucionais?

\section{Objetivos}

\subsection{Objetivo geral}

Analisar comparativamente o impacto das redes sociais na
participação política no Brasil e em Portugal.

\chapter{Referencial teórico}

O debate sobre participação política remonta às formulações
clássicas da teoria democrática, desde a democracia direta
ateniense até as concepções contemporâneas de democracia
deliberativa.

% --- Pós-textuais ---
\postextual

\chapter*{Referências}
\addcontentsline{toc}{chapter}{REFERÊNCIAS}

{\singlespacing
Referências seriam geradas aqui via \texttt{\textbackslash imprimirreferencias}
com um arquivo \texttt{.bib}.}

\end{document}
